\IfFileExists{stacks-project.cls}{%
\documentclass{stacks-project}
}{%
\documentclass{amsart}
}

% For dealing with references we use the comment environment
\usepackage{verbatim}
\newenvironment{reference}{\comment}{\endcomment}
%\newenvironment{reference}{}{}
\newenvironment{slogan}{\comment}{\endcomment}
\newenvironment{history}{\comment}{\endcomment}

% For commutative diagrams we use Xy-pic
\usepackage[all]{xy}

% We use 2cell for 2-commutative diagrams.
\xyoption{2cell}
\UseAllTwocells

% We use multicol for the list of chapters between chapters
\usepackage{multicol}

% This is generally recommended for better output
\usepackage{lmodern}
\usepackage[T1]{fontenc}

% For cross-file-references
\usepackage{xr-hyper}

% Package for hypertext links:
\usepackage{hyperref}

% For any local file, say "hello.tex" you want to link to please
% use \externaldocument[hello-]{hello}
\externaldocument[introduction-]{introduction}
\externaldocument[conventions-]{conventions}
\externaldocument[sets-]{sets}
\externaldocument[categories-]{categories}
\externaldocument[topology-]{topology}
\externaldocument[sheaves-]{sheaves}
\externaldocument[sites-]{sites}
\externaldocument[stacks-]{stacks}
\externaldocument[fields-]{fields}
\externaldocument[algebra-]{algebra}
\externaldocument[brauer-]{brauer}
\externaldocument[homology-]{homology}
\externaldocument[derived-]{derived}
\externaldocument[simplicial-]{simplicial}
\externaldocument[more-algebra-]{more-algebra}
\externaldocument[smoothing-]{smoothing}
\externaldocument[modules-]{modules}
\externaldocument[sites-modules-]{sites-modules}
\externaldocument[injectives-]{injectives}
\externaldocument[cohomology-]{cohomology}
\externaldocument[sites-cohomology-]{sites-cohomology}
\externaldocument[dga-]{dga}
\externaldocument[dpa-]{dpa}
\externaldocument[sdga-]{sdga}
\externaldocument[hypercovering-]{hypercovering}
\externaldocument[schemes-]{schemes}
\externaldocument[constructions-]{constructions}
\externaldocument[properties-]{properties}
\externaldocument[morphisms-]{morphisms}
\externaldocument[coherent-]{coherent}
\externaldocument[divisors-]{divisors}
\externaldocument[limits-]{limits}
\externaldocument[varieties-]{varieties}
\externaldocument[topologies-]{topologies}
\externaldocument[descent-]{descent}
\externaldocument[perfect-]{perfect}
\externaldocument[more-morphisms-]{more-morphisms}
\externaldocument[flat-]{flat}
\externaldocument[groupoids-]{groupoids}
\externaldocument[more-groupoids-]{more-groupoids}
\externaldocument[etale-]{etale}
\externaldocument[chow-]{chow}
\externaldocument[intersection-]{intersection}
\externaldocument[pic-]{pic}
\externaldocument[weil-]{weil}
\externaldocument[adequate-]{adequate}
\externaldocument[dualizing-]{dualizing}
\externaldocument[duality-]{duality}
\externaldocument[discriminant-]{discriminant}
\externaldocument[derham-]{derham}
\externaldocument[local-cohomology-]{local-cohomology}
\externaldocument[algebraization-]{algebraization}
\externaldocument[curves-]{curves}
\externaldocument[resolve-]{resolve}
\externaldocument[models-]{models}
\externaldocument[functors-]{functors}
\externaldocument[equiv-]{equiv}
\externaldocument[pione-]{pione}
\externaldocument[etale-cohomology-]{etale-cohomology}
\externaldocument[proetale-]{proetale}
\externaldocument[relative-cycles-]{relative-cycles}
\externaldocument[more-etale-]{more-etale}
\externaldocument[trace-]{trace}
\externaldocument[crystalline-]{crystalline}
\externaldocument[spaces-]{spaces}
\externaldocument[spaces-properties-]{spaces-properties}
\externaldocument[spaces-morphisms-]{spaces-morphisms}
\externaldocument[decent-spaces-]{decent-spaces}
\externaldocument[spaces-cohomology-]{spaces-cohomology}
\externaldocument[spaces-limits-]{spaces-limits}
\externaldocument[spaces-divisors-]{spaces-divisors}
\externaldocument[spaces-over-fields-]{spaces-over-fields}
\externaldocument[spaces-topologies-]{spaces-topologies}
\externaldocument[spaces-descent-]{spaces-descent}
\externaldocument[spaces-perfect-]{spaces-perfect}
\externaldocument[spaces-more-morphisms-]{spaces-more-morphisms}
\externaldocument[spaces-flat-]{spaces-flat}
\externaldocument[spaces-groupoids-]{spaces-groupoids}
\externaldocument[spaces-more-groupoids-]{spaces-more-groupoids}
\externaldocument[bootstrap-]{bootstrap}
\externaldocument[spaces-pushouts-]{spaces-pushouts}
\externaldocument[spaces-chow-]{spaces-chow}
\externaldocument[groupoids-quotients-]{groupoids-quotients}
\externaldocument[spaces-more-cohomology-]{spaces-more-cohomology}
\externaldocument[spaces-simplicial-]{spaces-simplicial}
\externaldocument[spaces-duality-]{spaces-duality}
\externaldocument[formal-spaces-]{formal-spaces}
\externaldocument[restricted-]{restricted}
\externaldocument[spaces-resolve-]{spaces-resolve}
\externaldocument[formal-defos-]{formal-defos}
\externaldocument[defos-]{defos}
\externaldocument[cotangent-]{cotangent}
\externaldocument[examples-defos-]{examples-defos}
\externaldocument[algebraic-]{algebraic}
\externaldocument[examples-stacks-]{examples-stacks}
\externaldocument[stacks-sheaves-]{stacks-sheaves}
\externaldocument[criteria-]{criteria}
\externaldocument[artin-]{artin}
\externaldocument[quot-]{quot}
\externaldocument[stacks-properties-]{stacks-properties}
\externaldocument[stacks-morphisms-]{stacks-morphisms}
\externaldocument[stacks-limits-]{stacks-limits}
\externaldocument[stacks-cohomology-]{stacks-cohomology}
\externaldocument[stacks-perfect-]{stacks-perfect}
\externaldocument[stacks-introduction-]{stacks-introduction}
\externaldocument[stacks-more-morphisms-]{stacks-more-morphisms}
\externaldocument[stacks-geometry-]{stacks-geometry}
\externaldocument[moduli-]{moduli}
\externaldocument[moduli-curves-]{moduli-curves}
\externaldocument[examples-]{examples}
\externaldocument[exercises-]{exercises}
\externaldocument[guide-]{guide}
\externaldocument[desirables-]{desirables}
\externaldocument[coding-]{coding}
\externaldocument[obsolete-]{obsolete}
\externaldocument[fdl-]{fdl}
\externaldocument[index-]{index}

% Theorem environments.
%
\theoremstyle{plain}
\newtheorem{theorem}[subsection]{Theorem}
\newtheorem{proposition}[subsection]{Proposition}
\newtheorem{lemma}[subsection]{Lemma}

\theoremstyle{definition}
\newtheorem{definition}[subsection]{Definition}
\newtheorem{example}[subsection]{Example}
\newtheorem{exercise}[subsection]{Exercise}
\newtheorem{situation}[subsection]{Situation}

\theoremstyle{remark}
\newtheorem{remark}[subsection]{Remark}
\newtheorem{remarks}[subsection]{Remarks}

\numberwithin{equation}{subsection}

% Macros
%
\def\lim{\mathop{\mathrm{lim}}\nolimits}
\def\colim{\mathop{\mathrm{colim}}\nolimits}
\def\Spec{\mathop{\mathrm{Spec}}}
\def\Hom{\mathop{\mathrm{Hom}}\nolimits}
\def\Ext{\mathop{\mathrm{Ext}}\nolimits}
\def\SheafHom{\mathop{\mathcal{H}\!\mathit{om}}\nolimits}
\def\SheafExt{\mathop{\mathcal{E}\!\mathit{xt}}\nolimits}
\def\Sch{\mathit{Sch}}
\def\Mor{\mathop{\mathrm{Mor}}\nolimits}
\def\Ob{\mathop{\mathrm{Ob}}\nolimits}
\def\Sh{\mathop{\mathit{Sh}}\nolimits}
\def\NL{\mathop{N\!L}\nolimits}
\def\CH{\mathop{\mathrm{CH}}\nolimits}
\def\proetale{{pro\text{-}\acute{e}tale}}
\def\etale{{\acute{e}tale}}
\def\QCoh{\mathit{QCoh}}
\def\Ker{\mathop{\mathrm{Ker}}}
\def\Im{\mathop{\mathrm{Im}}}
\def\Coker{\mathop{\mathrm{Coker}}}
\def\Coim{\mathop{\mathrm{Coim}}}

% Boxtimes
%
\DeclareMathSymbol{\boxtimes}{\mathbin}{AMSa}{"02}

%
% Macros for moduli stacks/spaces
%
\def\QCohstack{\mathcal{QC}\!\mathit{oh}}
\def\Cohstack{\mathcal{C}\!\mathit{oh}}
\def\Spacesstack{\mathcal{S}\!\mathit{paces}}
\def\Quotfunctor{\mathrm{Quot}}
\def\Hilbfunctor{\mathrm{Hilb}}
\def\Curvesstack{\mathcal{C}\!\mathit{urves}}
\def\Polarizedstack{\mathcal{P}\!\mathit{olarized}}
\def\Complexesstack{\mathcal{C}\!\mathit{omplexes}}
% \Pic is the operator that assigns to X its picard group, usage \Pic(X)
% \Picardstack_{X/B} denotes the Picard stack of X over B
% \Picardfunctor_{X/B} denotes the Picard functor of X over B
\def\Pic{\mathop{\mathrm{Pic}}\nolimits}
\def\Picardstack{\mathcal{P}\!\mathit{ic}}
\def\Picardfunctor{\mathrm{Pic}}
\def\Deformationcategory{\mathcal{D}\!\mathit{ef}}
\usepackage{cleveref}
\crefname{originalreference}{original reference}{original references}
\Crefname{originalreference}{Original reference}{Original references}
\makeatletter
\newlabel{situation-morphism-groupoids-on-field}{{04Q5}{source}{}{sourcepointer-1}{}}
\newlabel{situation-morphism-groupoids-on-field@cref}{{[originalreference][0][]04Q5}{[1][0][]source}{}{}{}}
\makeatother
\crefname{originallemma}{original lemma}{original lemmas}
\Crefname{originallemma}{Original lemma}{Original lemmas}
\makeatletter
\newlabel{lemma-groupoid-on-field-explain-points}{{04LR}{source}{}{sourcepointer-2}{}}
\newlabel{lemma-groupoid-on-field-explain-points@cref}{{[originallemma][0][]04LR}{[1][0][]source}{}{}{}}
\makeatother
\crefname{originallemma}{original lemma}{original lemmas}
\Crefname{originallemma}{Original lemma}{Original lemmas}
\makeatletter
\newlabel{lemma-restrict-groupoid-on-field}{{04LQ}{source}{}{sourcepointer-3}{}}
\newlabel{lemma-restrict-groupoid-on-field@cref}{{[originallemma][0][]04LQ}{[1][0][]source}{}{}{}}
\makeatother
\crefname{originallemma}{original lemma}{original lemmas}
\Crefname{originallemma}{Original lemma}{Original lemmas}
\makeatletter
\newlabel{topology-lemma-open-morphism-quotient-topology}{{02YB}{source}{}{sourcepointer-4}{}}
\newlabel{topology-lemma-open-morphism-quotient-topology@cref}{{[originallemma][0][]02YB}{[1][0][]source}{}{}{}}
\makeatother
\crefname{originallemma}{original lemma}{original lemmas}
\Crefname{originallemma}{Original lemma}{Original lemmas}
\makeatletter
\newlabel{lemma-groupoid-on-field-move-point}{{04LS}{source}{}{sourcepointer-5}{}}
\newlabel{lemma-groupoid-on-field-move-point@cref}{{[originallemma][0][]04LS}{[1][0][]source}{}{}{}}
\makeatother
\title{For a morphism of groupoids over a field as in the morphism-groupoids situation, an open image is closed}
\author{The Stacks Project Authors}
\date{stacks-project:more-groupoids.tex:local-snapshot-20261012; GNU FDL 1.2 or later}
\begin{document}
\maketitle
\setcounter{subsection}{1}
\noindent\textit{Editorial note.} All delivered mathematical text is the author's own, copied from the pinned Stacks Project chapter file more-groupoids.tex; no statement or proof step was rewritten, reordered or removed. The counted span is the author environment and its own proof environment. References inside the span are handed to the assembler as navigation stubs that carry the canonical tag of the cited tag only; whether each of them is load-bearing is recorded on the stub itself.\par\medskip
\noindent\textit{Allowed standard prerequisites (editorial).}\begin{itemize}
\item The morphism of groupoids over a field situation fixing R\_1, R\_2, a, s\_i, t\_i, c\_i (more-groupoids chapter, Situation situation-morphism-groupoids-on-field, canonical tag 04Q5; kind reference). cited at statement lines 1711-1712. Supplies the notation and assumptions in which the statement is phrased.
\item Base change to a field extension making source and target isomorphisms at an adapted point (more-groupoids chapter, Lemma lemma-groupoid-on-field-explain-points, canonical tag 04LR; collected on this lane as L22505). cited at proof lines 1718-1720. Provides the adapted pair k inside k' and r\_2' over r\_2 used to reduce to the case where s\_2 and t\_2 induce isomorphisms on residue fields.
\item Restriction of a groupoid over a field to a field extension has surjective open structure maps on the arrow scheme side (more-groupoids chapter, Lemma lemma-restrict-groupoid-on-field, canonical tag 04LQ). cited at proof lines 1732-1733. Makes p\_2 : R'\_2 to R\_2 surjective and open, so closure is detected after base change.
\item Open surjective morphisms and closure: preimages of closures are closures of preimages (topology chapter, Lemma topology-lemma-open-morphism-quotient-topology, canonical tag 02YB). cited at proof lines 1734-1735. Transfers r\_2 in the closure of a(R\_1) to r\_2' in the closure of a'(R'\_1) along the open surjective p\_2.
\item Automorphism moving the identity to a point when source and target are isomorphisms (more-groupoids chapter, Lemma lemma-groupoid-on-field-move-point, canonical tag 04LS; collected on this lane as L22506). cited at proof lines 1741-1743. Shows c(r\_2, a(R\_1)) is an open neighbourhood of r\_2, opening the way for the closure-intersection and composition argument.
\end{itemize}\medskip
\ifdefined\hypertarget\hypertarget{sourcepointer-1}{}\fi
\noindent\textit{Source reference.} Original reference 04Q5: sources/stacks-chow/tags.txt, source lines 6138--6138. Cited at statement lines 1711-1712: notation and assumptions fixing R\_1, R\_2, the morphism a and the structure maps.\par\medskip
\ifdefined\hypertarget\hypertarget{sourcepointer-2}{}\fi
\noindent\textit{Source reference.} Original lemma 04LR: sources/stacks-chow/tags.txt, source lines 6017--6017. Cited at proof lines 1718-1720: supplies the adapted extension k inside k' and the point r\_2' over r\_2 for the base-change reduction.\par\medskip
\ifdefined\hypertarget\hypertarget{sourcepointer-3}{}\fi
\noindent\textit{Source reference.} Original lemma 04LQ: sources/stacks-chow/tags.txt, source lines 6016--6016. Cited at proof lines 1732-1733: p\_2 : R'\_2 to R\_2 is surjective and open.\par\medskip
\ifdefined\hypertarget\hypertarget{sourcepointer-4}{}\fi
\noindent\textit{Source reference.} Original lemma 02YB: sources/stacks-chow/tags.txt, source lines 3870--3870. Cited at proof lines 1734-1735: r\_2 in the closure of a(R\_1) implies r\_2' in the closure of a'(R'\_1) along the open surjective p\_2.\par\medskip
\ifdefined\hypertarget\hypertarget{sourcepointer-5}{}\fi
\noindent\textit{Source reference.} Original lemma 04LS: sources/stacks-chow/tags.txt, source lines 6018--6018. Cited at proof lines 1741-1743: c(r\_2, a(R\_1)) is an open neighbourhood of r\_2 in the residue-field-isomorphism case.\par\medskip
\begin{lemma}
\label{lemma-open-image-is-closed}
Notation and assumptions as in
Situation \ref{situation-morphism-groupoids-on-field}.
If $a(R_1)$ is open in $R_2$, then $a(R_1)$ is closed in $R_2$.
\end{lemma}

\begin{proof}
Let $r_2 \in R_2$ be a point in the closure of $a(R_1)$.
We want to show $r_2 \in a(R_1)$. Pick $k \subset k'$ and
$r_2' \in R'_2$ adapted to $(U, R_2, s_2, t_2, c_2)$ and $r_2$ as in
Lemma \ref{lemma-groupoid-on-field-explain-points}.
Let $R_i'$ be the restriction of $R_i$ via the morphism
$U' = \Spec(k') \to U = \Spec(k)$.
Let $a' : R'_1 \to R_2'$ be the base change of $a$. The diagram
$$
\xymatrix{
R'_1 \ar[r]_{a'} \ar[d]_{p_1} & R'_2 \ar[d]^{p_2} \\
R_1 \ar[r]^a & R_2
}
$$
is a fibre square. Hence the image of $a'$ is the inverse image of
the image of $a$ via the morphism $p_2 : R'_2 \to R_2$. By
Lemma \ref{lemma-restrict-groupoid-on-field}
the map $p_2$ is surjective and open. Hence by
Topology, Lemma \ref{topology-lemma-open-morphism-quotient-topology}
we see that $r_2'$ is in the closure of $a'(R'_1)$.
This means that we may assume that $r_2 \in R_2$ has
the property that the maps $k \to \kappa(r_2)$ induced
by $s_2$ and $t_2$ are isomorphisms.

\medskip\noindent
In this case we can use
Lemma \ref{lemma-groupoid-on-field-move-point}.
This lemma implies $c(r_2, a(R_1))$ is an open neighbourhood of $r_2$.
Hence $a(R_1) \cap c(r_2, a(R_1)) \not = \emptyset$ as we assumed
that $r_2$ was a point of the closure of $a(R_1)$.
Using the inverse of $R_2$ and $R_1$ we see this means
$c_2(a(R_1), a(R_1))$ contains $r_2$.
As $c_2(a(R_1), a(R_1)) \subset a(c_1(R_1, R_1)) = a(R_1)$
we conclude $r_2 \in a(R_1)$ as desired.
\end{proof}

\end{document}
